\chapter{体验的计算几何：偏好、风险与历史依赖}
\label{chp:experience_geometry}
本章研究同一组客观可达路径为何会因主体历史、任务与风险预算而具有不同的选择难度。我们把世界图理解为对象和行动的可达结构，把体验图理解为在任务条件下改变比较、检索与行动优先级的内部状态；后者不是物理重力场，也不是痛苦在时空上留下的永久曲率。体验既包含瞬时感受或反馈，也包含由多次证据形成的慢偏好、风险模型与结构承诺。以下分别定义这些对象，说明先天约束与个体学习如何进入更新，再讨论有效距离、方向偏置、自我参照、逆向推断和工程验收。局部结构网络保存事件身份、来源和关系，全局分布表示承载相似性、候选与不确定性；两者通过任务相关绑定和读出接口联合。HSF--HD提供这种状态—结构共同演化的一般模型，AgentOS组件仅作为可执行实例。

\section{体验图的对象：瞬时反馈、慢偏好与任务度量}
\label{sec:section_9469}
我们先区分三个经常被“价值能量”混在一起的对象。瞬时反馈 $q_k\in\mathcal Q\subseteq\mathbb R^{d_q}$ 是时刻 $k$ 对刺激、内部稳态或环境回执的估计，它可以包含疼痛强度、愉悦报告、焦虑概率或人工系统中的失败类别；每一维都需要量表、方向和校准协议。慢偏好 $v_k\in\mathcal V\subseteq\mathbb R^{d_v}$ 汇总跨窗口证据，参与候选排序和风险分配。任务度量 $M_k\in\mathbb S_{++}^{d}$ 则定义在给定特征空间内怎样比较两种状态或轨迹，它是正定矩阵或其他满足所需公理的比较算子。三者属于不同空间，不能把瞬时强度直接称为度量曲率，也不能用设备焦耳能耗替代主观或算法反馈。

设对象状态为 $z_k\in\mathbb R^d$，任务版本为 $\theta_k$，候选动作集合为 $\mathcal A_k$，外部回执为 $y_k$。我们定义体验状态
\[
E_k=(q_k,v_k,M_k,\rho_k,\theta_k),
\]
其中 $\rho_k$ 是对各分量的来源和不确定性记录。局部反馈由观测模型 $q_k=Q_{\eta_k}(z_k,y_k,c_k)+\epsilon_k$ 产生，$c_k$ 是上下文，$\epsilon_k$ 是经声明的噪声。慢偏好不是对 $q_k$ 的无条件积分，而由带遗忘、来源权重和反事实校正的更新器生成：
\[
v_{k+1}=\Pi_{\mathcal V}\!\left[(1-\alpha_k)v_k+\alpha_k U_{\theta_k}(q_k,y_k,\rho_k)\right].
\]
这里 $\alpha_k\in[0,1]$ 是无量纲学习率，$\Pi_{\mathcal V}$ 保证偏好落在允许域。这个式子是模型假设；强反馈只有在来源可信、对象绑定正确并通过结构提交门槛时才产生长期改变。

历史文本把疼痛、快感与焦虑写成正交基向量。计算上，正交只表示选定内积下统计方向不相关，不能由名称自动得到。疼痛与焦虑可能高度相关，同一种疼痛在威胁学习、康复训练和医疗处置中也具有不同意义。我们因此先用校准数据估计协方差或任务度量，再决定是否需要正交化。坐标变换不删除经验，只改变表示；若变换使来源或对象身份丢失，即使数值上更简洁也不合格。对人工系统，失败代码、用户否决和环境损害同样不能被压成一个奖励标量后就视为完整体验，因为硬约束违例需要独立保留。

历史所谓价值联络可以改写为上下文相关的传送接口。不同语境中的偏好向量不能直接比较时，我们用 $P_{c\to c'}$ 把语境 $c$ 下的状态转换到 $c'$ 的坐标和权限体系，并以闭环残差检验转换是否一致。若一圈转换后产生显著残差，这只说明顺序、信息损失或接口不一致，不能未经群作用和曲率定义就称为非阿贝尔规范场。任务度量也随时间变化：当 $M_k$、$v_k$ 或 $\theta_k$ 更新时，轨迹代价的差分必须包含这些变化，不能只对 $z_k$ 求导后宣称价值保持。

工程验收包含四项。第一，用同一刺激的重复测量检查反馈校准，而不是只看输出幅度。第二，用对象交换和来源伪造测试绑定，避免把他人的事件写入自身历史。第三，用目标切换测试慢偏好是否能被任务版本正确解释，而不是永久污染所有场景。第四，用环境回执验证排序改变是否改善外部结果。局部结构网络在这里负责“谁经历了什么、何时发生、能否修改”，全局分布表示负责“哪些情境相似、哪些候选可能相关”；缺少任一方，体验图都容易退化为无身份的情绪向量或僵硬规则表。还要用相同平均反馈但不同时间顺序的轨迹测试历史依赖：若更新器对早期警告、近期恢复和持续低风险完全不加区分，慢偏好就不能支持情境化行动。对多主体系统，应分别保存报告者、承受后果者和决策者身份，防止群体平均掩盖个体损害。最终验收报告同时给出校准曲线、排序变化、约束违例和环境结果，不用一个总分替代相互冲突的证据。

\section{体验的形成：先验约束、证据更新与结构提交}
\label{sec:section_6973}
体验形成包含系统发生与个体发生两类来源，但二者都不等于质料向几何结构的物理相变。系统发生提供在部署或出生前已经存在的先验、反射、安全约束与可学习性结构；个体发生则在运行中依据事件证据修改预测、偏好和关系。我们把初始状态写为 $(v_0,M_0,B_0)$，其中 $B_0$ 是边界与硬约束。它们可来自遗传、设计、制度或预训练，不必都不可擦除。只有经过明确授权和风险评估的约束才被标记为硬约束，其余先验允许在反证下修订。自然选择能够解释某些生物倾向的历史来源，却不能证明每个个体对“死亡”都具有无穷势垒；人工系统也不能因为规则写在初始配置中就视为永恒真理。

个体学习采用候选—验证—提交三阶段。事件 $e_k=(I_k,o_k,a_k,y_k,t_k)$ 记录对象身份、观测、动作、环境结果和时间；快速更新先修改工作状态，慢速学习器再从窗口 $D_n=\{e_k\}_{k\in W_n}$ 生成候选参数和结构 $(\tilde v_{n+1},\tilde M_{n+1},\tilde G_{n+1})$。候选只有满足证据独立性、旧任务保持、新任务收益、边界安全与回滚可行五个条件才提交。连续参数可用投影更新，图编辑则通过版本化事务执行：
\[
(v_{n+1},M_{n+1},G_{n+1})=
\operatorname{Commit}\!\left((v_n,M_n,G_n),\mathcal P_n;\mathcal T_n\right),
\]
其中 $\mathcal P_n$ 是提案，$\mathcal T_n$ 是测试集。Commit是工程设计，不是普通加法；拒绝提案时保持原版本并保存失败证据。

“手碰火”保留为具体案例。一次接触产生伤害回执、疼痛报告和火源身份，快速控制立即撤回动作并提高短期风险估计。慢学习不应简单把所有“火”对象设为永久不可达，而要区分炉灶火焰、壁炉、远处图像和受控实验，学习观测条件、距离、保护装备与可行动作。若系统只记住红色或亮度，它会错误回避夕阳和图标；若只记住词语“火”，它又可能忽略无明火的高温表面。正确结构需要把对象、温度证据、接近动作、伤害结果和防护条件绑定起来，并用未见组合测试其泛化。

强烈事件也不自动获得最高长期权重。感知错误、恶意输入和重复转发都可能制造高幅反馈。更新器因此要估计来源相关性：来自同一记录的一百次复述不能算一百份独立证据；一次极端事件可以触发临时安全模式，却仍需复核后再改写长期结构。另一方面，平均化也可能抹去稀有灾难，故硬约束风险要按尾部概率或最坏情形单列。我们分别记录反馈强度、证据可信度、预期损失、计算资源与实测焦耳消耗，不以“高能”替代这些不同量。

所谓记忆疤痕在本模型中对应可追踪的路径依赖：相同当前输入因历史状态不同而产生不同估计或动作。我们通过配对运行测量这种差异，并检查它是否随安全证据逐步修订。若历史影响永不衰减，系统可能僵化；若一次相反反馈就完全擦除，系统又缺乏稳定性。学习率、遗忘率、证据门槛和边界模式共同决定稳定—可塑性权衡。变化环境下还要给旧经验设置适用域和有效期，避免在分布迁移后继续投射过时风险。

因此，体验形成是证据驱动的结构提交，而非能量把流形压出永久凹坑。协同论可以解释多个反馈通道如何形成稳定序参量，耗散与开放系统观点提醒这种维持需要信息和资源交换，认知科学则提供条件化、消退、泛化和重估的经验限制。HSF--HD把这些约束组织进状态—结构共同演化；AgentOS可用TCE实现候选更新、SPC保存结构提案、Boundary执行安全门槛、外部记忆保存来源，但其他架构也可实现同一协议。只有当干预能够定位历史证据如何改变读出，并且外部任务验证改善成立，我们才说经验真正进入了系统的计算几何。
